% ===========================================================================
% model.tex -- Section 2: model and preliminaries (G1).
% Extracted verbatim from the former model block of results.tex; the
% notation table lives in the appendix (Table~\ref{tab:notation}).
% ===========================================================================

\section{Model and preliminaries}\label{sec:model}

A ground set $N$ with $|N|=n$, a monotone submodular $f:2^{N}\to\mathbb
R_{\ge0}$ with $f(\emptyset)=0$, and a cardinality budget $K$ are fixed.  The
algorithm cannot evaluate $f$; it can only query a predictor
$\tilde f:2^{N}\to\mathbb R$ with $\tilde f(\emptyset)=0$.  Marginal gains are
written $d_e(S)=f(S\cup\{e\})-f(S)$ and
$\tilde d_e(S)=\tilde f(S\cup\{e\})-\tilde f(S)$.
% Decision D1 (TASKS5): the main model places no structural assumption on
% ftilde; the submodular-surrogate variant appears only in
% Remark \ref{rem:exact-gap}.

\begin{definition}[Prediction error]\label{def:eta}
For $\eta_u,\eta_o\ge1$ the predictor $\tilde f$ has (single-element,
multiplicative) error at most $(\eta_u,\eta_o)$ if for every $S\subseteq N$
and $e\notin S$,
\[
  d_e(S)/\eta_u\;\le\;\tilde d_e(S)\;\le\;\eta_o\,d_e(S).
\]
In particular $d_e(S)=0$ forces $\tilde d_e(S)=0$, and all predicted gains are
nonnegative.  The scalar error is $\eta=\eta_u\eta_o$.
\end{definition}
The multiplicative form of Definition~\ref{def:eta} follows the error
conventions of learning-augmented scheduling, where predicted weights or
processing times are accurate up to multiplicative factors and guarantees
degrade with that factor
\citep{lattanzi2020online,azar2021flow,zhao2022real}.
% H1: spirit-of citations restored from the TPAMI submission's error-metric
% section (it cited exactly these three at this point); all three re-passed
% the F6 four-step check (results/H1_bib_crosscheck.md).

Approximation ratios are stated as $\alpha\in(0,1]$ with
$F^{\mathrm{ALG}}\ge\alpha\,F^{\mathrm{OPT}}$; larger is better.
% Erratum guard: the submitted draft stated the ratio in the opposite
% direction in Section 1.1; fixed here (RESEARCH_STATE, known errors).

Predictive greedy is the single-step greedy run on $\tilde f$: for
$t=0,\dots,K-1$ it adds an element maximizing $\tilde d_e(S^{t})$, with ties
broken adversarially in all worst-case statements.  Its exact worst-case ratio
at error level $\eta$ is written $\rho_K(\eta)$.  Two run-dependent error
measures refine $\eta$:
\begin{definition}[Selection error]\label{def:etasel}
For an execution of predictive greedy with states
$S^{0},\dots,S^{K-1}$ and picks $e_{0},\dots,e_{K-1}$, write
$M_t=\max_{e\notin S^{t}} d_{e}(S^{t})$ and $g_t=d_{e_t}(S^{t})$
(monotonicity of $f$ gives $g_t\ge0$), and set
\[
  a_t\;=\;\begin{cases}
    M_t/g_t, & g_t>0,\\
    1, & M_t=g_t=0,\\
    \infty, & g_t=0<M_t .
  \end{cases}
  \qquad
  \etasel\;=\;\max\{1,a_0,\dots,a_{K-1}\},
\]
the largest factor by which a step of the run fell short of the best true
gain; a step whose chosen true gain vanishes while a better candidate
exists makes the selection error infinite.  Bounds evaluated at
$\etasel$ use the convention $L_K(\infty)=0$.  In the terminology of
Goundan and Schulz, a run with finite $\etasel$ is an execution of greedy
with an $\alpha$-approximate incremental oracle with $\alpha=\etasel$;
their condition is a per-step condition, which is what the case
$a_t=\infty$ records \citep{goundan2007revisiting}.
% Status of the attribution: F6 audit item 4 and the J2 external review
% both confirm alpha = etasel on GS p. 7 (not 1/etasel).  The notation
% \etasel is ours.  Their Theorem 1 is the L_K bound; Proposition
% \ref{prop:guarantee} is that theorem instantiated and is attributed
% accordingly, not claimed as new (decision D2).
% Decision D3 (J2 adoption, 2026-09-07): the earlier definition excluded
% the steps with g_t = 0 and its unconditional certificate reading was
% refuted by a three-element modular instance (J2 review section 4.1,
% re-verified in results/H3_j2_recheck.py); the stepwise a_t close that
% hole with the infinity case.
\end{definition}
The trajectory error restricts Definition~\ref{def:eta} to the states of the
run; it is used in a remark and in the appendix, and written $\etatr$ when a
symbol is needed.  For every run of a predictor with finite error
$(\eta_u,\eta_o)$ the chain $\etasel\le\etatr\le\eta$ holds; in particular
$a_t=\infty$ cannot occur then, because a vanishing chosen predicted gain
forces every candidate's predicted, and then true, gain to vanish.
% [HAND-PROOF-UNREVIEWED, three lines: for finite steps max_e d_e <=
%  eta_u max_e dtilde_e <= eta_u dtilde_{e_t} <= eta_u eta_o d_{e_t} at
%  every state of the run; for g_t = 0 the band gives M_t = 0, so
%  a_t = 1.  Without a finite band the chain can fail (J2 section 4).]

Table~\ref{tab:notation} in Appendix~\ref{sec:notation} collects the
notation.
